\section{Conclusion}
\label{sec:conclusion}

The architecture described in this paper rests on a small number of design commitments. First, applications --- not the platform --- define the algebra of their state; the platform is generic over a commutative monoid supplied as WebAssembly. Second, synchronization is parameterized by that same algebra through a contract-supplied summary/delta protocol, so the wire cost of reconciliation tracks application semantics rather than raw state size. Third, the network underneath is a small-world ring whose long-range links arise from ordinary request traffic and whose forwarding decisions are informed by a per-neighbor performance model rather than topological distance alone. Fourth, private state and sensitive operations live in actor-style delegates that hold secrets behind a platform-enforced trust boundary, leaving public state in contracts.

Taken together these commitments support a class of decentralized applications --- group chat, collaborative documents, reputation systems, social feeds --- whose consistency requirements have always been awkward to express either on top of content-addressed immutability or underneath blockchain linearizability. The platform's goal is not to be the right substrate for every decentralized application; it is to make the broad middle of that space buildable by ordinary developers in ordinary languages, with the heavy lifting --- routing, replication, sync, transport, sandboxing --- carried out beneath them.

Several substantial questions remain open. Quantitative characterization of the accept-only-at-terminus topology, formal verification of contract algebras, and a workable incentive layer that resists coordinated abuse are the directions that most plainly bound what the platform can become. We expect to address those questions empirically, in successor work, on top of the architecture described here.
